% !TeX root = main.tex


% William 
\section{Introduction}

% Intro to SW and HPC SW
Improving all aspects of scientific software development has become essential for the sustainability and trustworthiness of modern science~\cite{osti_1846008}. Early on, software
development for high-performance computing (HPC) was identified as a non-trivial task that requires a deep understanding of the
application and the targeted architectural systems~\cite{parashar1994study}. The complexity of HPC software is expected to increase
with the hardware heterogeneity characterizing the present systems~\cite{osti_1473756}. Therefore, understanding this
landscape~\cite{grannan2020understanding} is essential to guarantee the effective and strategic use of industry-standard software engineering
practices that have a positive impact on a particular HPC community target~\cite{basili2008understanding}. 

QMCPACK\footnote{\url{https://github.com/QMCPACK/qmcpack}}~\cite{kim2018qmcpack,doi:10.1063/5.0004860} is an open-source many-body
ab-initio Quantum Monte Carlo~\cite{foulkes2001} (QMC) framework. QMC methods solve directly the Schrodinger equation offering more accuracy than density functional theory (DFT) methods at a greater computational cost. QMCPACK is used to compute the electronic structure of atoms, molecules, 2D nanomaterials and solids.
The code supports many architectures ranging from laptops through mid-range clusters to the largest supercomputers in the world. Written in C\texttt{++} with a modular design, it uses a hybrid parallel approach combining the message passing interface (MPI) with multicore, many-core CPUs and GPUs. It primarily uses OpenMP~\cite{chapman2008using} for CPU and GPU execution, with additional CUDA~\cite{nickolls2008scalable},
HIP/ROCm~\cite{hip} and SYCL~\cite{reyes2016sycl} for higher GPU performance. Releases are typically made quarterly. QMCPACK is part of the application development portfolio~\cite{osti_1594850} of the US Department of Energy Exascale Computing Project (ECP)~\cite{8528398}. ECP's mission 
addresses currently intractable problems of strategic importance and national interest at the same time as the first exascale systems, such as Frontier\footnote{\url{https://www.olcf.ornl.gov/frontier}}, are deployed.

Cloud technologies have contributed greatly to the development and empowerment of the open-source community in the last decade. They provide a rich ecosystem and resources that facilitate the application of modern software engineering practices to enable the trustworthiness that is crucial in scientific code~\cite{fox2011cloud}. Code hosting services like GitHub have evolved into a one-stop shop that adds value to the modern software development process. One example is the rapid adoption of GitHub Actions~\cite{decan2022use} since its initial release in 2019, providing computational resources for continuous
integration/continuous development (CI/CD) workflows. It also provides seamless integration with third-party cloud services
(\textit{e.g.}, code analysis, coverage, reporting) via a marketplace of reusable workflow ``actions''.

In this paper, we describe our efforts and the impact of incorporating modern software engineering practices in the QMCPACK code. The goal of these changes has been to
improve the quality of the software and enable significant refactoring while also keeping the barrier for open source software contributions low.
Metrics are provided to quantify the value of such practices in a real HPC application with active users and developers. We showcase: (i) the evolution
and failure rates on our GitHub Actions-based CI system to cover CPU and GPU test cases, (ii) the code sanitization process to
resolve and future-proof for memory leaks, (iii) the inclusion of Docker containers to enable reproducible virtual testing
environments, (iv) code refactoring efforts to ensure the sustainability of the scientific efforts around the software. As a result of these improvements, the QMCPACK community benefits from a robust first line of quality checks ahead of deployment, testing, and science production on HPC systems.

The paper is organized as follows: Section~\ref{sec:Background} provides background information on the QMCPACK software.
Section~\ref{sec:Software-Improvements} describes the software improvements showcased in this paper along with metrics for each
covered aspect. Related efforts in applying software engineering practices to HPC codes are presented in Section~\ref{sec:Related-Work}. Section~\ref{sec:Conclusions} presents the conclusions of this work, while an artifact description
for the activities and reproducibility of the metrics reported in this paper is given in Appendix~\ref{ap1:Artifact}. 
